\documentclass[11pt, a4paper]{article}

% --- Encoding, Sprache & Schriftart ---
\usepackage[utf8]{inputenc}
\usepackage[T1]{fontenc}
\usepackage[english]{babel}
\usepackage{mathptmx} % Serifen-Schriftart (Times-ähnlich) wie im Original-Template
\usepackage[scaled=0.92]{helvet} % Sans-Serif für Titelseite

% --- Seitenränder & Layout ---
\usepackage[top=2.5cm, bottom=2.5cm, left=2.5cm, right=2.5cm]{geometry}
\usepackage{parskip}
\usepackage{titlesec}
\usepackage{fancyhdr}
\usepackage{tabularx}
\usepackage{booktabs}
\usepackage{array}
\usepackage{graphicx}
\usepackage{xcolor}
\usepackage{hyperref}

% --- Farben & Links ---
\definecolor{linkblue}{RGB}{0, 51, 102}
\definecolor{reqgray}{RGB}{245, 246, 248}
\definecolor{bordergray}{RGB}{180, 185, 190}

\hypersetup{
    colorlinks=true,
    linkcolor=black,
    urlcolor=linkblue,
    pdftitle={Software Requirements Specification},
    pdfauthor={Team Name}
}

% --- Projekt-Metadaten (Hier einmalig anpassen!) ---
\newcommand{\projecttitle}{<Project Name>}
\newcommand{\docversion}{Version 1.0 draft}
\newcommand{\docauthor}{<Team Name / Authors>}
\newcommand{\docorg}{DHBW Ravensburg -- Campus Friedrichshafen}
\newcommand{\docdate}{\today}

% --- Kopf- und Fußzeilen (wie auf Seiten ii bis 6 des Templates) ---
\pagestyle{fancy}
\fancyhf{}
\renewcommand{\headrulewidth}{0pt}
\fancyhead[L]{\textit{\textbf{Software Requirements Specification for \projecttitle}}}
\fancyhead[R]{\textit{\textbf{Page \thepage}}}

% --- Überschriften-Formatierung (exakt wie im IEEE-Template) ---
\titleformat{\section}
  {\normalfont\Large\bfseries}{\thesection.}{0.8em}{}
\titleformat{\subsection}
  {\normalfont\large\bfseries}{\thesubsection}{0.8em}{}
\titleformat{\subsubsection}
  {\normalfont\normalsize\bfseries}{\thesubsubsection}{0.8em}{}

% --- CUSTOM MAKROS FÜR REQUIREMENTS MANAGEMENT ---
% 1. Anforderungs-Makro mit automatischem Sprungziel (Label)
% Nutzung: \req{FR-01}{User Login}{High}{The system shall authenticate users via DHBW SSO.}
\newcommand{\req}[4]{%
    \par\vspace{4pt}%
    \noindent\fcolorbox{bordergray}{reqgray}{%
        \parbox{\dimexpr\textwidth-2\fboxsep-2\fboxrule\relax}{%
            \phantomsection\label{req:#1}%
            \textbf{#1: #2} \hfill \textit{[Priority: #3]}\\[3pt]
            #4%
        }%
    }\par\vspace{4pt}%
}

% 2. Verweis auf ein Requirement im Fließtext (klickbar)
% Nutzung: \reqref{FR-01}
\newcommand{\reqref}[1]{\hyperref[req:#1]{\textbf{#1}}}

% 3. Platzhalter-Hinweis (kursiv wie im Original-Template, zum Löschen beim Ausfüllen)
\newcommand{\guidance}[1]{\textit{\color{darkgray}<#1>}}

\begin{document}

% ==============================================================================
% TITELSEITE (Exakt nach Page 1 des IEEE-Templates)
% ==============================================================================
\begin{titlepage}
    \sffamily
    \noindent\rule{\textwidth}{3pt}
    \vspace{1.2cm}
    
    \begin{flushright}
        {\fontsize{28}{34}\selectfont\bfseries Software Requirements\\Specification\par}
        \vspace{1.2cm}
        {\Large\bfseries for\par}
        \vspace{1.2cm}
        {\fontsize{24}{30}\selectfont\bfseries \projecttitle\par}
        \vspace{1.8cm}
        {\large\bfseries \docversion\par}
        \vspace{1.2cm}
        {\large\bfseries Prepared by \docauthor\par}
        \vspace{1.2cm}
        {\large\bfseries \docorg\par}
        \vspace{1.2cm}
        {\large\bfseries \docdate\par}
    \end{flushright}
    
    \vfill
    \rmfamily
    \begin{center}
        \small\textit{\textbf{Copyright \copyright\ 1999 by Karl E. Wiegers. Permission is granted to use, modify, and distribute this document.}}
    \end{center}
\end{titlepage}

% ==============================================================================
% INHALTSVERZEICHNIS & REVISION HISTORY (Römische Seitenzahlen, Page ii)
% ==============================================================================
\pagenumbering{roman}
\setcounter{page}{2} % Startet bei "Page ii" wie im Original

\renewcommand{\contentsname}{Table of Contents}
\tableofcontents

\vspace{1.2cm}
\section*{Revision History}
\addcontentsline{toc}{section}{Revision History}

\begin{tabularx}{\textwidth}{|p{3.2cm}|p{2.4cm}|X|p{2.2cm}|}
    \hline
    \textbf{Name} & \textbf{Date} & \textbf{Reason For Changes} & \textbf{Version} \\
    \hline\hline
    <Author Name> & 2026-11-17 & Initial draft created from inception document & 0.1 \\
    \hline
     & & & \\
    \hline
     & & & \\
    \hline
\end{tabularx}

\newpage

% ==============================================================================
% HAUPTTEIL (Arabische Seitenzahlen ab Page 1)
% ==============================================================================
\pagenumbering{arabic}
\setcounter{page}{1}

% ------------------------------------------------------------------------------
\section{Introduction}
% ------------------------------------------------------------------------------

\subsection{Purpose}
\guidance{Identify the product whose software requirements are specified in this document, including the revision or release number. Describe the scope of the product that is covered by this SRS.}

\subsection{Document Conventions}
\guidance{Describe any standards or typographical conventions that were followed when writing this SRS. Example: Unique IDs are assigned using FR-XX for Functional Requirements and NFR-XX for Non-Functional Requirements. Priorities follow the MoSCoW or High/Medium/Low scheme.}

\subsection{Intended Audience and Reading Suggestions}
\guidance{Describe the different types of reader that the document is intended for, such as developers, project managers, users, testers, and documentation writers.}

\subsection{Product Scope}
\guidance{Provide a short description of the software being specified and its purpose, including relevant benefits, objectives, and goals.}

\subsection{References}
\guidance{List any other documents or Web addresses to which this SRS refers (e.g., Inception Document, DHBW Project Guidelines, IEEE Std 830-1998).}
\begin{itemize}
    \item Inception Document for \projecttitle, Version 1.0, 2026.
    \item Prof. Dr. Christian Braunagel: \textit{Software Engineering -- Program Design Part I}, DHBW Ravensburg, Q4 2026.
\end{itemize}

% ------------------------------------------------------------------------------
\section{Overall Description}
% ------------------------------------------------------------------------------

\subsection{Product Perspective}
\guidance{Describe the context and origin of the product being specified in this SRS. Include a high-level context diagram if helpful.}

\subsection{Product Functions}
\guidance{Summarize the major functions the product must perform or must let the user perform. Only a high-level summary (such as a bullet list) is needed here.}
\begin{itemize}
    \item Function 1: ...
    \item Function 2: ...
\end{itemize}

\subsection{User Classes and Characteristics}
\guidance{Identify the various user classes that you anticipate will use this product (e.g., Students, Lecturers, Administrators) and describe their pertinent characteristics.}

\subsection{Operating Environment}
\guidance{Describe the environment in which the software will operate, including hardware platform, operating system and versions, browsers, or databases.}

\subsection{Design and Implementation Constraints}
\guidance{Describe any items or issues that will limit the options available to the developers (e.g., DHBW IT policies, mandatory frameworks, time constraints for Part II).}

\subsection{User Documentation}
\guidance{List the user documentation components (such as user manuals, online help, or READMEs) that will be delivered along with the software.}

\subsection{Assumptions and Dependencies}
\guidance{List any assumed factors (as opposed to known facts) that could affect the requirements stated in the SRS, as well as dependencies on external systems (e.g., DHBW Moodle API or LDAP).}

% ------------------------------------------------------------------------------
\section{External Interface Requirements}
% ------------------------------------------------------------------------------

\subsection{User Interfaces}
\guidance{Describe the logical characteristics of each interface between the software product and the users. You can reference UI mockups here.}

\subsection{Hardware Interfaces}
\guidance{Describe the logical and physical characteristics of each interface between the software product and the hardware components of the system.}

\subsection{Software Interfaces}
\guidance{Describe the connections between this product and other specific software components (databases, operating systems, APIs, libraries).}

\subsection{Communications Interfaces}
\guidance{Describe the requirements associated with any communications functions required by this product (e.g., HTTPS, REST APIs, WebSockets, email notifications).}

% ------------------------------------------------------------------------------
\section{System Features}
% ------------------------------------------------------------------------------
\guidance{Organize the functional requirements for the product by system features or use cases. Copy subsection 4.1 for each additional feature.}

\subsection{User Authentication and Profile Management}
\subsubsection{Description and Priority}
This feature allows DHBW students and lecturers to securely log into the platform using their university credentials and manage their profile settings. \textbf{Priority: High} (Benefit: 9, Penalty: 8, Cost: 3, Risk: 2).

\subsubsection{Stimulus/Response Sequences}
\begin{enumerate}
    \item \textbf{Stimulus:} User clicks ``Login with DHBW Account'' on the landing page.
    \item \textbf{Response:} System redirects to the university authentication provider.
    \item \textbf{Stimulus:} User enters valid credentials and confirms.
    \item \textbf{Response:} System creates a user session and displays the personal dashboard.
\end{enumerate}

\subsubsection{Functional Requirements}
% Beispiel für die Verwendung des Requirements-Makros:
\req{FR-01}{DHBW Single Sign-On}{High}{
    The system shall allow users to authenticate using their official DHBW university credentials. If authentication fails, an informative error message shall be displayed without revealing sensitive account details.
}

\req{FR-02}{Role Assignment}{High}{
    Upon first successful login, the system shall automatically assign the user role (Student or Lecturer) based on the directory attributes returned by the identity provider (see \reqref{FR-01}).
}

\subsection{System Feature 2 (Replace with Feature Name)}
\subsubsection{Description and Priority}
\guidance{Provide a short description of the feature and indicate whether it is of High, Medium, or Low priority.}

\subsubsection{Stimulus/Response Sequences}
\guidance{List the sequences of user actions and system responses that stimulate the behavior defined for this feature.}

\subsubsection{Functional Requirements}
\req{FR-03}{Example Requirement Title}{Medium}{
    The system shall ...
}

% ------------------------------------------------------------------------------
\section{Other Nonfunctional Requirements}
% ------------------------------------------------------------------------------

\subsection{Performance Requirements}
\req{NFR-01}{Page Load Response Time}{Medium}{
    95\% of all user interactions and page requests shall complete within 1.5 seconds under a simulated load of 100 concurrent campus users.
}

\subsection{Safety Requirements}
\guidance{Specify those requirements that are concerned with possible loss, damage, or harm that could result from the use of the product.}

\subsection{Security Requirements}
\req{NFR-02}{Transport Encryption}{High}{
    All network communication between client browsers and the backend server shall be encrypted using TLS 1.3 or higher.
}

\subsection{Software Quality Attributes}
\guidance{Specify additional quality characteristics such as adaptability, availability, maintainability, reliability, and usability.}

\subsection{Business Rules}
\guidance{List any operating principles about the product, such as which individuals or roles can perform which functions under specific circumstances.}

% ------------------------------------------------------------------------------
\section{Other Requirements}
% ------------------------------------------------------------------------------
\guidance{Define any other requirements not covered elsewhere in the SRS (e.g., database persistence requirements, GDPR/DSGVO compliance, internationalization).}

% ------------------------------------------------------------------------------
% ANHÄNGE (Appendices A, B, C)
% ------------------------------------------------------------------------------
\newpage
\appendix
\renewcommand{\thesection}{Appendix \Alph{section}:}
\titleformat{\section}{\normalfont\Large\bfseries}{\thesection}{0.5em}{}

\section{Glossary}
\guidance{Define all the terms necessary to properly interpret the SRS, including acronyms and abbreviations.}
\begin{description}
    \item[DHBW] Duale Hochschule Baden-Württemberg.
    \item[SRS] Software Requirements Specification.
    \item[TBD] To Be Determined.
\end{description}

\section{Analysis Models}
\guidance{Include any pertinent analysis models, such as use-case diagrams, data flow diagrams, class diagrams, or entity-relationship diagrams.}
% Beispiel zum Einbinden eines Diagramms aus draw.io:
% \begin{figure}[h!]
%     \centering
%     \includegraphics[width=0.85\textwidth]{use_case_diagram.png}
%     \caption{Initial Use-Case Diagram}
%     \label{fig:usecase}
% \end{figure}

\section{To Be Determined List}
\guidance{Collect a numbered list of the TBD references that remain in the SRS so they can be tracked to closure.}
\begin{enumerate}
    \item \textbf{TBD-01:} Clarify with customer team whether offline caching is required for mobile devices (refers to Section 2.4).
\end{enumerate}

\end{document}
```
### Vollständig KI generiertes Template mittels Gemini basiert auf dem gegebenen Dokument
### So arbeitet ihr damit optimal in inscrive:
2. **Requirements verwalten & verlinken:** Mit dem Befehl `\req{ID}{Titel}{Priorität}{Text}` wird jede Anforderung in einem einheitlichen Rahmen dargestellt und automatisch als Sprungziel registriert. Überall im Dokument könnt ihr dann einfach `\reqref{FR-01}` schreiben, und im PDF entsteht ein klickbarer Verweis.
3. **Hilfstexte entfernen:** Alle kursiven `<...>`-Hinweise aus der Originalvorlage sind über das Makro `\guidance{...}` eingebunden – ersetzt sie einfach schrittweise durch eure eigenen Texte.